\section{API-design}
\label{sec:api-design}

API-et er designet som et REST-basert grensesnitt der videorom er hovedressursen, i tråd med kravspesifikasjonens beskrivelse av operasjoner for oppretting, sletting, henting, start og stopp av videorom \cite{hdo_kravspec}. Alle operasjoner er samlet under \texttt{/api/rooms}, slik at klienten forholder seg til én enhetlig kontrakt uavhengig av underliggende videomotor.

% API-et er designet som et REST-basert grensesnitt der videorom er hovedressursen.
% Dette passer med kravspesifikasjonen, som beskriver operasjoner for å opprette,
% slette, hente informasjon om, starte og stoppe videorom~\cite{hdo_kravspec}.
% Operasjonene er derfor samlet under \texttt{/api/rooms}, slik at klienten får ett
% felles grensesnitt uavhengig av hvilken videomotor som brukes.

ASP.NET Core støtter både Minimal APIs og controller-baserte API-er. Minimal APIs
er anbefalt for nye prosjekter der målet er raske HTTP-endepunkter med lite
oppsett og mindre kode, mens controller-baserte API-er gir en mer strukturert
modell med egne controller-klasser, attributtbasert routing og innebygd støtte
for API-spesifikke funksjoner som modellbinding, validering og standardiserte
feilsvar~\cite{microsoft_aspnet_apis,microsoft_aspnet_webapi}. I VideoAPI ble
controller-basert API valgt fordi løsningen har flere ressursorienterte
endepunkter, felles autentisering på controller-nivå og behov for tydelig
separasjon mellom HTTP-kontrakt og provider-logikk. Minimal APIs kunne vært brukt
for en mindre prototype, men controller-modellen ga en tydeligere struktur for
operasjoner som oppretting, sletting, start, stopp, tokenhåndtering og
WebRTC-lenker for videorom.

VideoAPI eksponerer utelukkende et server-side REST API over HTTP uten
klientbibliotek. Primære konsumenter er backend-tjenester og operatørapplikasjoner;
enhver HTTP-klient kan integrere mot det. AntMedia og NHN er de to videomotorene
som er implementert i proof-of-concept-et, men provider-arkitekturen er designet
for at nye motorer kan legges til uten endringer i API-kontrakten
(jf.\ avsnitt~\ref{sec:provider-validering}).

Et sentralt designvalg var hvordan API-et skulle bevare informasjon om hvilken backend et rom tilhører. Løsningen ble å returnere rom-ID-er på formatet \texttt{\{provider\}:\{backendId\}}, for eksempel \texttt{antmedia:abc123} eller \texttt{nhn:24484}. Provider-prefikset fungerer som rutinginformasjon i senere forespørsler, slik at klienten kun trenger å lagre én ID. Kodeliste~\ref{lst:create-room} viser at controlleren holdes tynn: den binder forespørselen, velger provider, delegerer romopprettelsen og prefikser rom-ID-en før svaret returneres.
\begin{lstlisting}[language=CSharp, caption={Opprettelse av rom — controlleren delegerer til provider-laget.}, label={lst:create-room}]
// HDO-VideoAPI/Controllers/RoomsController.cs

[HttpPost]
public async Task<IActionResult> CreateRoom([FromBody] CreateRoomRequest request)
{
    // ... (genere navn om ikke oppgitt)

    // Velg provider basert på forespørsel, standardverdi er antmedia
    var providerKey = (request.Provider ?? "antmedia").ToLowerInvariant();
    var provider = factory.GetProvider(providerKey);

    // Deleger romoppretting til provider-laget
    var room = await provider.CreateRoomAsync(request, HttpContext.RequestAborted);

    // Prefiks ID med provider-nøkkel
    room.Id = $"{providerKey}:{room.Id}";

    // ... (logging)

    // Returnere 201 Created med Location-header
    return CreatedAtAction(nameof(GetRoom), new { roomId = room.Id }, room);
}
\end{lstlisting}

Ved senere kall — henting, sletting, start og stopp — må controlleren gjenkjenne hvilken provider rom-ID-en tilhører. Kodeliste~\ref{lst:parse-room-id} viser hvordan provider-prefikset parses ut av ID-en, slik at provider-valget ikke trenger å sendes som et eget felt eller query-parameter.

\begin{lstlisting}[language=CSharp, caption={Parsing av provider-prefiks fra rom-ID}, label={lst:parse-room-id}]
// HDO-VideoAPI/Controllers/RoomsController.cs

private (IVideoProvider Provider, string ProviderKey, string RawId) ParseRoomId(string roomId)
{
    // Finn skilletegnet mellom provider og backend-ID
    var colon = roomId.IndexOf(':');

    // ... (valider at kolon finnes)

    var providerKey = roomId[..colon];  // f.eks. "antmedia"
    var rawId = roomId[(colon + 1)..];  // f.eks. "abc123"

    // ... (valider at rawId ikke er tom)

    // Slå opp provider og returner alle tre delene
    return (factory.GetProvider(providerKey), providerKey, rawId);
}
\end{lstlisting}

Start og stopp av videostrøm er modellert som handlinger på romressursen, for
eksempel \texttt{POST /api/rooms/\{roomId\}/start} og
\texttt{POST /api/rooms/\{roomId\}/stop}. Dette ble valgt fordi operasjonene
alltid gjelder et bestemt rom. Samtidig gir ikke alle operasjoner mening for alle
backends. AntMedia har en broadcast-modell der start og stopp passer naturlig,
mens NHN Booking API i større grad beskriver et videomøte som er klart etter
opprettelse. Slike forskjeller håndteres i provider-laget.

API-et bruker JSON for forespørsler og svar, og responsmodellene er normalisert
så langt det er hensiktsmessig. Klienten får samme overordnede struktur
uavhengig av backend, men backend-spesifikke forskjeller skjules ikke dersom de
påvirker hvordan operasjonen må brukes. Dette gjelder særlig WebRTC-lenker:
AntMedia returnerer en kortlivet tokenbasert tilgang til streamen, mens NHN
returnerer en møte-URL med tilgangsinformasjon.

Feilresponsene er også standardisert. Dersom en provider ikke støtter en
operasjon, eller en ekstern backend returnerer feil, oversettes dette til et
strukturert HTTP-svar gjennom den globale feilhåndteringen. Dette gjør at
controllerne kan være like for alle providere, mens klienten fortsatt får
forutsigbare HTTP-statuskoder og feilmeldinger.